\chapter{Compilação, revisão e erros}
\label{chap:erros}

\section{Ler a primeira mensagem útil}

Quando uma compilação falha, o terminal pode exibir muito texto. Procure a
primeira linha iniciada por \verb|!| ou uma referência no formato
\arquivo{arquivo.tex:linha}. Erros posteriores frequentemente são apenas
consequências do primeiro.

Problemas comuns incluem:

\begin{itemize}
  \item \emph{Undefined control sequence}: comando digitado incorretamente ou
    pacote ausente;
  \item \emph{File not found}: caminho, extensão ou maiúsculas não coincidem;
  \item \emph{Missing \}}: chave aberta sem fechamento;
  \item \emph{Environment ... undefined}: ambiente incorreto ou pacote não
    carregado;
  \item \emph{There were undefined references}: rótulo ausente ou necessidade
    de outra passagem.
\end{itemize}

Reduza o problema: comente temporariamente partes recentes, compile e reative
pequenos blocos. Não apague arquivos-fonte para ``limpar'' o projeto. Use
\arquivo{latexmk -C} para remover apenas artefatos.

\section{Avisos que merecem revisão}

Uma compilação bem-sucedida ainda pode conter avisos. \emph{Overfull hbox}
indica conteúdo ultrapassando a margem, muitas vezes por URL ou código longo.
Referências e citações indefinidas aparecem como interrogações no PDF. Fontes
substituídas podem alterar a aparência esperada.

Revise também o resultado visual: páginas isoladas, títulos no fim da página,
imagens ilegíveis, inconsistência de termos e espaços excessivos. A ausência de
erros técnicos não garante um bom livro.

\section{Um fluxo diário seguro}

Trabalhe em incrementos pequenos:

\begin{enumerate}
  \item escreva uma seção curta;
  \item compile com \arquivo{latexmk -lualatex main.tex};
  \item corrija o primeiro erro e repita;
  \item confira avisos e páginas alteradas;
  \item registre uma versão funcional no Git.
\end{enumerate}

O modo contínuo, \arquivo{latexmk -lualatex -pvc main.tex}, recompila ao salvar.
No VS Code, a extensão LaTeX Workshop oferece compilação, visualização e salto
entre fonte e PDF.

\begin{atencao}
  Um projeto que depende de LuaLaTeX deve ser compilado sempre com LuaLaTeX.
  Executar pdfLaTeX sobre um preâmbulo com \arquivo{fontspec} produz um erro por
  projeto de forma intencional.
\end{atencao}

\begin{pratica}
  Faça uma cópia de um capítulo e provoque três erros: remova uma chave, altere
  o nome de um ambiente e referencie uma imagem inexistente. Identifique a
  primeira mensagem útil de cada caso e depois restaure o arquivo.
\end{pratica}

